%%%%%%%%%%  scan idx 24  |  folio 11  %%%%%%%%%%
\SGAPrintPage{11}
\[
T_f\in K^\bullet(X).
\]
When $f$ admits a factorization (2.1) through an immersion into a smooth relative scheme, the definition of $T_f$ presents no new difficulty. We can moreover reformulate it in this case, in a manner more in keeping with the ``derived categories'' point of view, by considering the familiar canonical homomorphism
\sgaStableTarget{sga6.e00.eq.4-1}\begin{equation}\tag{4.1}
\underline N_{X/X'}\xrightarrow{d}\Omega^1_{X'/Y}\otimes_{\underline O_{X'}}\underline O_X
\end{equation}
induced by the exterior differential
\[
d_{X'/Y}:\underline J\subset\underline O_{X'}\longrightarrow\Omega^1_{X'/Y},
\]
and the chain complex
\sgaStableTarget{sga6.e00.opening.eq.4-2}\begin{equation}\tag{4.2}
L^{X/Y}_\bullet\in\operatorname{Ob}D(X)
\end{equation}
defined from (4.1) ``by extending by zeros'', a complex equal in degree $1$, respectively $0$, to the source, respectively the target, of the homomorphism (4.1). Formula (2.2) then takes the more pleasant form
\sgaStableTarget{sga6.e00.opening.eq.4-3}\begin{equation}\tag{4.3}
\check T_f=\cl^\bullet\bigl(L^{X/Y}_\bullet\bigr),
\end{equation}
a formula which makes sense in $K^\bullet(X)$ when $f$ is a locally complete intersection morphism, since in this case the complex $L^{X/Y}_\bullet$ is manifestly perfect. The assertion of invariance of $T_f$ relative to the chosen factorization of $f$ can then readily be made more precise as an assertion of invariance of the complex $L^{X/Y}_\bullet$: up to canonical isomorphism in the derived category $D(X)$, this complex is likewise independent of the chosen factorization (Exp. VIII).

\sgaStableTarget{sga6.e00.sec.4.4}\subsubsection*{4.4.}
When $f:X\to Y$ is an arbitrary morphism of schemes, we can still construct a canonical object $L^{X/Y}_\bullet$ of the derived category $D(X)$, defined up to unique isomorphism, which on each affine open of $X$ over an affine open of $Y$ is isomorphic to the complex constructed according to the principle explained in the preceding paragraph (where, however, in the general case ``smooth'' should be replaced by ``formally smooth'', so that every algebra $B$ over a ring $A$ occurs as a quotient of a formally smooth algebra, for example a polynomial algebra). For such a definition, we cannot proceed by simple ``gluing'' from the affine constructions, because the objects of the derived category $D(X)$ are by their nature

%%%%%%%%%%  scan idx 25  |  folio 12  %%%%%%%%%%
\SGAPrintPage{12}
essentially not glueable. Here is a general construction of $L^{X/Y}_\bullet$, valid more generally for every morphism $f:X\to Y$ of (commutatively) ringed topoi: let $B=\underline O_X$, $A=f^{-1}(\underline O_Y)$, so that $A$ is a Ring on the space $X$, and $B$ is an $A$-Algebra. Let $T\to B$ be a homomorphism of sheaves of \underline{sets} generating $B$ as an $A$-Algebra, i.e. such that the corresponding homomorphism of $A$-Algebras
\[
\phi:C=A[T]\longrightarrow B
\]
is an epimorphism of sheaves of sets, where $A[T]$ denotes the free commutative $A$-Algebra generated by $T$ (i.e. the sheaf associated with the presheaf $U\longmapsto A(U)[T(U)]$ = free commutative $A(U)$-algebra generated by $T(U)$).
Let $J$ be the kernel of $\phi$, and consider, as in 4.3, the chain complex of length $1$ defined by the natural homomorphism
\[
J/J^2\xrightarrow{d}\Omega^1_{C/A}\otimes_C B.
\]
It is not difficult to see that, up to canonical isomorphism in the category $D(X)$, the complex so constructed is independent of the choice of $T$ and of $T\to B$, so that it deserves the name \uline{relative cotangent complex $L^{X/Y}_\bullet$ of $X$ over $Y$}. We also verify that, on every affine open of $X$ over an affine open of $Y$, it is given (up to canonical isomorphism) by the procedure of 4.3.

When $f$ is locally complete intersection, we moreover see that we even obtain a \underline{perfect} complex, which then allows us again to define an element $T_f$ in $K^\bullet(X)$ by formula (4.3). We point out that here again, in general, it would not have been possible to define $T_f$ as an element of $K^\bullet_{\mathrm{naive}}(X)$, and nothing permits us to suppose that $T_f$ is \underline{globally} isomorphic in $D(X)$ to a bounded complex with locally free components of finite type. This therefore provides further justification for the point of view advanced in 4.1 in favour of the ``sophisticated'' $K^\bullet(X)$, in preference to $K^\bullet_{\mathrm{naive}}(X)$.

\sgaStableTarget{sga6.e00.sec.4.5}\subsubsection*{4.5.}
Subject to the question raised in 4.2, we have therefore again assembled all the structural elements needed to give a meaning to (1.2) for a proper locally complete intersection morphism of arbitrary schemes (the noetherian hypothesis implicitly made in the preceding sections can in fact be removed without essential difficulty). Unfortunately, this does not mean that, even subject to this reservation, the formula in question has been established, even when $Y$ is the spectrum of a field

%%%%%%%%%%  scan idx 26  |  folio 13  %%%%%%%%%%
\SGAPrintPage{13}
that is algebraically closed of characteristic $p>0$, and $X$ is, say, a proper smooth algebraic scheme of dimension $3$. See the comments in Exp. XIV no.~2.

\sgaStableTarget{sga6.e00.sec.5}\subsection*{5.}
In the form presented in no.~4, the Riemann--Roch formula (1.2) also has a meaning for a proper locally complete intersection morphism of complex analytic spaces (or of rigid analytic spaces in the sense of Tate). But of course the proof given in [RRR] applies to this case only when we also assume that the morphism $f$ is \underline{projective}, and, as in the algebraic case, it seems that an essentially new idea is needed to treat the general case. The formula we envisage here has values in an ``analytic Chow ring'' $\Gr^\bullet(X)$, or rather $\Gr^\bullet(X)_{\mathbb Q}$, constructed according to the principle explained in 4.2 (the difficulty encountered there moreover disappears because we are in characteristic zero, as noted in Exp. XIV no.~1).

When $X$ and $Y$ are nonsingular, another version of the Riemann--Roch formula can be given as follows. First let $X$ be an arbitrary complex analytic space, and $X^{\mathrm{top}}$ the same space equipped with the sheaf of continuous complex-valued functions. We thus have a canonical morphism of ringed spaces
\[
X^{\mathrm{top}}\longrightarrow X,
\]
which, by the evident contravariance of the functor $K^\bullet(X)$ in the variable ringed space $X$, defines a ring homomorphism
\sgaStableTarget{sga6.e00.eq.5-1}\begin{equation}\tag{5.1}
K^\bullet(X)\longrightarrow K^\bullet(X^{\mathrm{top}}).
\end{equation}
For simplicity, suppose that $X^{\mathrm{top}}$ is compact; then, as in the case (already mentioned in 3.1) of a scheme admitting an ample invertible module, we see (Exp. II) that every perfect complex on $X^{\mathrm{top}}$ is globally isomorphic to a bounded complex locally free in each degree, i.e. to a complex of complex vector bundles on $X^{\mathrm{top}}$, and consequently $K^\bullet(X^{\mathrm{top}})$ is canonically isomorphic to the topologists' familiar group $K(X^{\mathrm{top}})$ [1], defined in terms of complex vector bundles on the compact space $X^{\mathrm{top}}$. Thus we can interpret the homomorphism (5.1) as a homomorphism from $K^\bullet(X)$, defined in terms of the analytic structure, to the topologists' familiar ring $K(X^{\mathrm{top}})$. Now, for a morphism $f:X\longrightarrow Y$ of compact nonsingular analytic spaces, we know how to associate, by Atiyah--Hirzebruch

%%%%%%%%%%  scan idx 27  |  folio 14  %%%%%%%%%%
\SGAPrintPage{14}
[1], a group homomorphism
\sgaStableTarget{sga6.e00.eq.5-2}\begin{equation}\tag{5.2}
f^{\mathrm{top}}_*:K^\bullet(X^{\mathrm{top}})\longrightarrow K^\bullet(Y^{\mathrm{top}}).
\end{equation}
Using likewise the homomorphism $f_*$ of (3.1) defined by formula (3.2) (which here implicitly uses Grauert's finiteness theorem [3] for a proper morphism of analytic spaces), we obtain a square of natural homomorphisms involving only groups of type $K^\bullet$ (to the exclusion of rings of Chow-ring or cohomology type):
\sgaStableTarget{sga6.e00.eq.5-3}\begin{equation}\tag{5.3}
\begin{array}{ccc}
K^\bullet(X) & \longrightarrow & K^\bullet(X^{\mathrm{top}}) \\
f_*\downarrow && \downarrow f^{\mathrm{top}}_* \\
K^\bullet(Y) & \longrightarrow & K^\bullet(Y^{\mathrm{top}}).
\end{array}
\end{equation}
The Riemann--Roch type formula to which we alluded consists in asserting that this square commutes. We note the difference in nature between this assertion (which does not disregard torsion phenomena, and lies in a group $K^\bullet(Y^{\mathrm{top}})$ of an essentially discrete nature) and the Riemann--Roch formula of the type considered in no.~4, which lies in a group of ``Chow'' type that is no longer discrete in nature but must instead be tensored with $\mathbb Q$. We note that formula (5.3) is proved, in any case, when $X$ and $Y$ are nonsingular projective varieties [7], as is the formula of the type considered in no.~4, which is proved in this case in [RRR]. This is a plausible variant of the Atiyah--Singer index theorem (and of its generalizations due to Shih [6] and Atiyah), which plainly calls for a common generalization. Once the commutativity of (5.3) has been proved, the ``differentiable Riemann--Roch theorem'' of Atiyah--Hirzebruch [1] would immediately give a cohomological variant of the complex analytic Riemann--Roch formula, namely the commutativity of the square
\sgaStableTarget{sga6.e00.eq.5-4}\begin{equation}\tag{5.4}
\begin{array}{ccc}
K^\bullet(X) & \xrightarrow{\;x\longrightarrow \ch(x)\Todd(T_X)\;} & H^{2*}(X,\mathbb Q) \\
\downarrow && \downarrow \\
K^\bullet(Y) & \xrightarrow{\;y\longrightarrow \ch(y)\Todd(T_Y)\;} & H^{2*}(Y,\mathbb Q),
\end{array}
\end{equation}
where the horizontal arrows are obtained from the cohomological Chern classes, obtained from the homomorphism (5.1) and the ordinary Chern class homomorphism
\[
K^\bullet(X^{\mathrm{top}})\longrightarrow H^{2*}(X,\mathbb Z),
\]

%%%%%%%%%%  scan idx 28  |  folio 15  %%%%%%%%%%
\SGAPrintPage{15}
the first vertical arrow is the same as in (5.3), and the second is the Gysin homomorphism in rational cohomology.

\sgaStableTarget{sga6.e00.sec.6}\subsection*{6.}
Throughout this discussion, we have considered almost exclusively the contravariant invariant $K^\bullet(X)$, to the exclusion of the covariant invariant $K_\bullet(X)$ whose existence was mentioned in 3.1. Regarding $K^\bullet(X)$ as an analogue of the integral \underline{cohomology} ring of a topological space, we must regard $K_\bullet(X)$ as an analogue of integral \underline{homology}, the Module structure of $K_\bullet(X)$ over $K^\bullet(X)$ then corresponding to the cap-product operation\SGAAuthorFootnote{raisedstar}{This analogy is moreover not purely verbal: after passing to the associated graded objects $\Gr^\bullet(X)$, $\Gr_\bullet(X)$, it can be made precise by an explicit compatibility relation between the module structure on $\Gr_\bullet(X)$ and the ``cap'' operation.}. The fact that, when $X$ is regular, the natural homomorphism $K^\bullet(X)\to K_\bullet(X)$ is an isomorphism must be regarded as the analogue of the Poincaré duality theorem on oriented topological manifolds, asserting that cap product with the fundamental homology class defines an isomorphism from cohomology to homology. These analogies also suggest the desirability of defining, say for a finite polyhedron $X$, an invariant $K_\bullet(X)$ of a covariant nature, whose relations to integral homology (spectral sequence, Chern homomorphism\dots) would be analogous to those between $K^\bullet(X)$ and integral cohomology; nor does it seem excluded that there might exist, in terms of these groups $K_\bullet(X)$, variants of the differentiable Riemann--Roch theorem valid for spaces more general than compact differentiable manifolds.

Returning to an arbitrary noetherian scheme $X$, we should equip $K_\bullet(X)$ with an increasing filtration, $\operatorname{Filt}_iK_\bullet(X)$ being generated by classes $\cl_\bullet(F)$ with $F$ a coherent sheaf on $X$ such that $\dim \operatorname{supp}.F\le i$. We shall prove (Exp. X) the compatibility of this filtration with the module structure and the filtration of $K^\bullet(X)$, i.e. the formula
\sgaStableTarget{sga6.e00.eq.6-1}\begin{equation}\tag{6.1}
\operatorname{Filt}^iK^\bullet(X)\cdot \operatorname{Filt}_j K_\bullet(X)\subset \operatorname{Filt}_{j-i}K_\bullet(X),
\end{equation}
which equips the associated graded object of the filtered module $K_\bullet(X)$, denoted by
\[
\Gr_\bullet(X),
\]
with a module structure over $\Gr^\bullet(X)$. Its elements can moreover be interpreted as \uline{classes of algebraic cycles} on $X$, (for an equivalence relation less fine than rational equivalence, when $X$ is of finite type over a base field $K$). If

